
\section{Introduction to Capstone}
\label{capstone}

\textbf{Quick Reference Guide}

\begin{center} 
\begin{tabular}{|>{\raggedright}p{1.8in}|>{\raggedright}p{2.5in}|c|} 
\hline \centering{\textbf{What You Want To Do}} & \centering{\textbf{How You Do It}} & \textbf{Button} \\ 

\hline Start Recording Data 
& Click the `Record' button
& \raisebox{-5mm}{\includegraphics{../../131/StudentGuideModule1/appendices/capstone/record_button.eps}} \\ 

\hline Stop Recording Data 
& Click the `Stop' button
& \raisebox{-5mm}{\includegraphics{../../131/StudentGuideModule1/appendices/capstone/stop_button.eps}} \\ 

\hline Re-scale the data so it fills the Graph display window 
& Click the `Scale to Fit' button
& \raisebox{-5mm}{\includegraphics{../../131/StudentGuideModule1/appendices/capstone/scale_to_fit.eps}} \\ 

\hline Pinpoint the x- and y-coordinate values on the Graph display
& Click the `Delta Tool' button, select ``Add Coordinates/Delta Tool'', move square to point of interest
& \raisebox{-5mm}{\includegraphics{../../131/StudentGuideModule1/appendices/capstone/delta_tool.eps}} \\ 

\hline Select a section of data
& Click the `Highlight' button. See instructions below.
& \raisebox{-5mm}{\includegraphics{../../131/StudentGuideModule1/appendices/capstone/highlight.eps}} \\ 

\hline Select from the Statistics menu
& Click the `Statistics' button, use menu. See instructions below.
& \raisebox{-5mm}{\includegraphics{../../131/StudentGuideModule1/appendices/capstone/statistics.eps}} \\

\hline Find area under a curve
& Click the `Display Area' button after selecting data of interest
& \raisebox{-5mm}{\includegraphics{../../131/StudentGuideModule1/appendices/capstone/display_area.eps}} \\



\hline 
\end{tabular} 
\end{center}
\index{color page}

\begin{center}
To delete Delta Tool: Hold cursor on Delta Tool square, right click and select ``Delete Tool''.
\end{center}

\vspace{5mm}
\textbf{Selecting a Section of Data}

\begin{enumerate}
\item To select a data section, click the \textbf{Highlight} button.  Drag the pink rectangle to the data you wish to highlight.  Adjust the width and height using horizontal and vertical arrows at edges of highlighted region.
\item To unselect the data, with cursor in highlighted region, right click and select ``Delete highlighter''.
\end{enumerate}
\vspace{5mm}
\textbf{Finding Average and Standard Deviation of a Section of Data}

\begin{enumerate}
\item Select the section of data of interest, as described above.
\item Click on the \textbf{Statistics} button on the Graph Toolbar. Use the menu on this button to select maximum, minimum, mean, standard deviation, or any combination.
\item To remove the results, click on the \textbf{Statistics} button again.
\end{enumerate}


\pagebreak[2]
\vspace{5mm}
\textbf{Finding the Area Under a Curve}

\begin{enumerate}
\item Select the section of data that you want to integrate under, using the procedure described above.
\item Click the \textbf{Display Area} button on the Graph Toolbar.
The results of the integration will be displayed on the graph.
\item To undo the integration, click on the \textbf{Display Area} button again.
\end{enumerate}

\pagebreak[2]

\textbf{Calibrating Force Sensors}

\begin{enumerate}
\item Check to see that the force sensor is connected to the Pasco interface.

\item Open whatever experiment file you will be using.

\item Click the ``Calibration'' icon in the tools panel on the left side of the screen.

\item Verify that the type of measurement is set to ``Force'', then click ``Next''.

\item Choose ``Two standards'' as calibration type, then click ``Next''.

\item In this step, you will apply known forces to the sensors using weights.  You can either hang masses from the sensor directly, or using a string looped over a pulley.  If you are hanging masses directly while holding the force sensor in your hand, try resting it against the edge of the table to keep it stable.

\begin{enumerate}

\item With no mass attached, set the standard value to 0~N for the first point, click ``set current value to standard value'', then click ``Next''.

\item Hang a mass (say 200 g) from the sensor, set the standard value to its weight (say 1.96~N), click ``set current value to standard value'', then click ``Next''.

\end{enumerate}

\item After reviewing calibration, click ``Finish''.

\item Click the ``Calibration'' icon again to close the dialog. 

\item If you calibrated the sensor by holding it vertically but will now use it horizontally, you will need to adjust for the weight of the hook.  Turn the sensor to horizontal, without any applied force, and hit the ``TARE'' button on the side of the force sensor.

\item The force sensor should now calibrated for the rest of your experiment.  But it's a good idea to try taking a tiny bit of data with a known weight to make sure your results aren't totally nuts.

\end{enumerate}
